% chapters/ch05_uart_rs_serial.tex — Chapter 5: UART / RS232 / RS422 / RS485
% Scope (PLAN.md): frame anatomy, baud math & error %, parity/framing/overrun
% errors, flow control, single-ended vs differential, RS422 vs RS485 vs RS232,
% why flight computers use RS422 (her daily interface), termination & noise.
% Budget 35 pp, mix 22 Q&A / 10 code / 3 concept.
% All snippets verified gcc -Wall -Wextra -std=c11 (zero warnings).

\chapter{UART / RS232 / RS422 / RS485}
\label{ch:uart-serial}

\begin{partnote}
\textbf{Part B opens here.} \textbf{Scope} --- asynchronous frame anatomy
$\rightarrow$ baud-rate arithmetic and error percentage $\rightarrow$
parity/framing/overrun errors $\rightarrow$ flow control $\rightarrow$
single-ended vs differential signalling $\rightarrow$ RS232 vs RS422 vs RS485
$\rightarrow$ why flight computers use RS422 $\rightarrow$ termination and
noise. This is the candidate's \emph{daily} interface --- RS422 Embedded C on
the DFCC at TASL --- so it must be bulletproof.
\end{partnote}

The terms get blurred in casual use, so fix them first: \textbf{UART} is the
\emph{controller} (the on-chip peripheral and the asynchronous framing);
\textbf{RS232/RS422/RS485} are \emph{electrical standards} (voltage levels and
single-ended-vs-differential signalling) that a line driver applies to that
same framing. A UART speaks bytes; a transceiver decides the volts on the wire.
Conflating them is the first thing an interviewer probes.

%=====================================================================
\section{Primer}
%=====================================================================

\subsection{The asynchronous frame}

UART is \emph{asynchronous}: there is no shared clock, so each side is
pre-configured to the same \textbf{baud rate} and agrees on a frame format.
The line idles high. A character is sent as: one \textbf{start} bit (a
high$\rightarrow$low edge that the receiver locks onto), then the
\textbf{data} bits \emph{least-significant-bit first} (5--9, usually 8), an
optional \textbf{parity} bit, then one or more \textbf{stop} bits (line
returns high). The shorthand ``8N1'' means 8 data bits, No parity, 1 stop bit
--- 10 bit-times per byte once the start and stop are counted. The receiver
samples each bit near its centre (typically at the 8th of 16 oversampling
ticks), so both ends must hold the same baud closely enough that the sampling
point hasn't drifted by the last bit.

\subsection{Baud, error, and why it must match}

Baud rate is bits per second on the wire (for UART, symbol rate = bit rate).
The peripheral derives it by dividing the clock: on an STM32, the 16$\times$
oversampling divider \code{USARTDIV} satisfies
$\text{baud}=f_{\text{ck}}/(16\cdot\text{USARTDIV})$, and the register
\code{BRR}$=f_{\text{ck}}/\text{baud}$ to the nearest integer. Integer rounding
means the \emph{actual} baud differs slightly from nominal, and that error
accumulates across the frame: by the stop bit the sampling point has shifted by
(error)$\times$(bits). The rule of thumb is the cumulative error must stay
under about half a bit, so each end's baud error should be well under a few
percent --- which is why an accurate \textbf{crystal} (Chapter 3 \code{HSE}),
not the internal RC, is used for reliable UART.

\subsection{Electrical standards: single-ended vs differential}

\textbf{RS232} is single-ended: one wire per direction, referenced to ground,
swinging roughly $\pm$3--15\,V (idle/mark is negative). Simple, short-range
(a few metres), one driver and one receiver. \textbf{RS422} and \textbf{RS485}
are \emph{differential}: each signal is a twisted pair (A/B), and the receiver
reads the \emph{difference} between the two lines. Common-mode noise picked up
on the cable hits both wires equally and cancels at the differential receiver,
so differential links run far (over a kilometre) and fast with strong noise
immunity --- the reason avionics and industrial buses use them. \textbf{RS422}
is point-to-point / multi-drop with one driver and up to ten receivers
(full-duplex, separate TX and RX pairs). \textbf{RS485} is multi-point: up to
32 (or more) drivers \emph{and} receivers share one pair (half-duplex), so
drivers must be tri-stated and turn-taking is managed in software. This is why
flight computers favour RS422: differential noise immunity, deterministic
point-to-point links, and full-duplex separation of command and telemetry.

%=====================================================================
\section{Q\&A Bank}
%=====================================================================

\tierbanner{Tier 1 --- Screening (phone / HR-technical)}%
{Two to four sentences. Serial basics --- expected cold from anyone claiming
embedded comms.}

\question{Q1.1}{What is the difference between a UART and RS232/RS422/RS485?}
The UART is the digital peripheral that frames bytes asynchronously
(start/data/parity/stop) at TTL levels; RS232/RS422/RS485 are electrical
standards that a line driver/transceiver uses to put those bits on a cable at
defined voltages and (for 422/485) differentially. Same framing, different
physical layer.

\question{Q1.2}{Describe a standard UART frame.}
Idle-high line, then a start bit (falling edge), 5--9 data bits sent
LSB-first, an optional parity bit, and 1--2 stop bits returning the line high.
``8N1'' is 8 data bits, no parity, 1 stop --- 10 bit-times per byte.

\question{Q1.3}{What does ``baud rate'' mean and why must both ends agree?}
It's the number of signalling bits per second on the wire. Because UART is
asynchronous with no shared clock, each side times its bits from its own clock
at the agreed baud; if they differ too much, the receiver's sampling point
drifts off the bit and data corrupts.

\question{Q1.4}{Single-ended vs differential signalling --- one line.}
Single-ended (RS232) measures one wire against ground; differential
(RS422/485) measures the difference between two wires, so common-mode noise
cancels --- giving far longer range and far better noise immunity.

\question{Q1.5}{RS422 vs RS485 in one breath?}
RS422 is differential, typically point-to-point/multi-drop, full-duplex (one
driver, separate TX/RX pairs). RS485 is differential multi-point, half-duplex
(many drivers share one pair), so drivers must be enabled only when
transmitting.

\question{Q1.6}{What is parity and what can it catch?}
An extra bit set so the total number of 1s is even (even parity) or odd (odd
parity). It detects any single-bit (odd number of) error in the frame but
cannot detect two-bit errors or correct anything --- a weak, cheap check.

\question{Q1.7}{Name three UART receive errors.}
Framing error (stop bit sampled as 0 --- baud mismatch or noise), parity error
(parity bit doesn't match), and overrun (a new byte arrived before the previous
one was read from the data register).

\question{Q1.8}{Why do RS422/RS485 cables use twisted pairs?}
Twisting makes both wires pick up virtually identical interference, so the
noise appears as common-mode and cancels at the differential receiver. It also
reduces the loop area that radiates/receives EMI.

\question{Q1.9}{What is flow control and name the two common kinds?}
A mechanism to stop the sender overrunning the receiver. Hardware flow control
uses RTS/CTS lines; software flow control uses XON/XOFF characters in the data
stream. Without it, a slow receiver drops bytes (overrun).

\question{Q1.10}{Why is RS232 limited to short distances?}
Single-ended signalling references ground, so ground shifts and noise add
directly to the signal, and its higher-voltage swing has limited slew at
distance --- practically a few metres at moderate baud, versus over a kilometre
for differential RS422/485.

\tierbanner{Tier 2 --- Core technical (panel round)}%
{A short paragraph. Mechanisms, arithmetic, and the bring-up reality.}

\question{Q2.1}{Work the baud divider: 16\,MHz clock, 115200 baud, 16$\times$
oversampling --- what register value and what's the real baud and error?}
\code{BRR}$=f_{\text{ck}}/\text{baud}=16{,}000{,}000/115200=138.9$, rounded to
\textbf{139}. The actual baud is $16{,}000{,}000/139 = 115{,}107$, an error of
about $-807$\,ppm ($-0.08\%$). That's tiny --- well under the ~2\% budget --- so
8N1 works fine. The lesson: pick a clock that divides your baud cleanly (e.g.\
a 14.7456\,MHz crystal hits standard bauds exactly), and always check the error
at high baud where the divisor is small and rounding hurts most.

\question{Q2.2}{How much baud mismatch can a UART tolerate, and why?}
The receiver samples each bit near its centre, so the accumulated timing error
at the \emph{last} sampled bit must stay within about half a bit. With a start
bit plus 8 data plus stop ($\approx$10 bit-times) and centre sampling, the
combined error of both ends should stay under roughly $\pm$2--3\%. Beyond that
the sample point walks off the bit and you get framing/parity errors. Two ends
each at 1.5\% in opposite directions is already marginal --- another reason to
crystal-clock both sides.

\question{Q2.3}{Walk through computing and checking even parity in firmware.}
Parity is the XOR-reduction of the data bits: fold the byte
(\code{b\^{}=b>>4; b\^{}=b>>2; b\^{}=b>>1; p=b\&1}), giving 1 when the number of
set bits is odd --- which \emph{is} the even-parity bit (it makes the total
even). Odd parity is the complement. On receive you recompute and compare to
the received parity bit; a mismatch flags a parity error. It's $O(1)$ and
branch-free, but remember it only catches odd numbers of bit-flips.

\question{Q2.4}{What is a framing error and what usually causes it?}
The UART expects the stop bit to be high; if it samples low there, it raises a
framing error. Causes: baud-rate mismatch (the sample point drifted so the
``stop'' lands on a data bit), line noise/glitches, a break condition (line
held low), or wrong frame format (expecting 1 stop bit when 2 were sent).
On bring-up a steady stream of framing errors almost always means the baud or
clock config is wrong, not the cable.

\question{Q2.5}{What is an overrun error and how do you prevent it?}
Overrun happens when a newly received byte overwrites one the CPU hadn't read
from the data register yet --- data is lost. Prevention: service RX promptly
via interrupt (push to a ring buffer, Chapter 4), or use DMA to move bytes
without per-byte CPU work, and apply flow control so a busy receiver can pause
the sender. At high baud on a slow core, interrupt-per-byte latency is the
usual culprit --- DMA is the fix.

\question{Q2.6}{Hardware (RTS/CTS) vs software (XON/XOFF) flow control ---
trade-offs?}
RTS/CTS uses dedicated lines: the receiver de-asserts RTS/CTS to tell the
sender ``pause,'' so it's fast, in-band-data-clean, and reacts within a byte ---
but needs extra wires. XON/XOFF embeds two control bytes in the data stream:
no extra wires, but it steals two byte values from the data (bad for binary/
telemetry), and reaction is delayed by buffering. For binary avionics
telemetry you use hardware flow control or a higher-layer protocol, never
XON/XOFF.

\question{Q2.7}{Why does RS485 need the driver-enable (DE) managed, and what's
the timing risk?}
RS485 is half-duplex: many drivers share one pair, so each must tri-state
(release the bus) except while transmitting, controlled by the transceiver's DE
pin. The risk is \emph{turnaround}: you must assert DE before the first bit and
hold it until the last stop bit has fully shifted out (plus a guard time)
before releasing --- drop DE early and you truncate the last byte; release late
and you collide with the next talker. The guard time is a few bit-periods,
computed from the baud (Coding Gym).

\question{Q2.8}{How does 9-bit / multidrop addressing work on RS485?}
A 9th data bit marks a byte as an \emph{address} rather than \emph{data}:
the master sends an address byte with the 9th bit set; all slaves' UARTs (in
mute/address-match mode) wake and compare it, and only the addressed slave
listens to the following data bytes (9th bit clear). It lets many nodes share
one bus without every node interrupting on every byte --- hardware
``mute until my address'' support makes it efficient.

\question{Q2.9}{Why do differential lines need termination, and where?}
A fast edge travelling down a transmission line reflects off an unmatched end,
and the reflection corrupts the signal. Terminating each end of the pair with a
resistor equal to the cable's characteristic impedance (typically 120\,$\Omega$
for twisted pair) absorbs the energy and kills reflections. RS422 terminates at
the receiver end of each pair; RS485 terminates at \emph{both} physical ends of
the bus (not at stubs in the middle). Missing or wrong termination shows up as
ringing on the scope and intermittent errors at high baud or long cable.

\question{Q2.10}{Why do flight computers favour RS422 over RS232 for inter-card
links?}
Differential signalling gives strong common-mode noise rejection (vital in the
electrically harsh, EMI-heavy environment of an aircraft --- MIL-STD-461),
long reliable range, higher usable baud, and full-duplex separation of command
and telemetry on independent pairs. RS232's single-ended, ground-referenced,
short-range nature is fine for a debug console but not for a flight-critical
data link. This is exactly why the DFCC card-to-card interface on the resume is
RS422.

\tierbanner{Tier 3 --- Trap \& depth (principal-engineer level)}%
{A paragraph plus the trap. The defensible answers behind ``RS422 daily work.''}

\question{Q3.1}{RS422 and RS485 use the same differential transceivers and
voltages --- so what \emph{actually} distinguishes them, and when does it bite?}
Electrically they're nearly identical (differential pairs, $\sim\pm$2--6\,V);
the real differences are \emph{topology and duplex}. RS422 is one driver per
pair (point-to-point or one-to-many \emph{receivers}), usually full-duplex with
separate TX/RX pairs, so the driver is always enabled and there's no bus
contention. RS485 is multi-\emph{driver} on a shared pair, half-duplex, so it
\emph{requires} driver-enable arbitration, turnaround timing, and bus
termination at both ends. \trap the trap is treating them as interchangeable
because the transceiver datasheet looks similar: wire an RS485 multidrop like an
RS422 point-to-point (drivers always enabled) and two nodes transmitting at
once short the bus. The distinguishing design work is the \emph{protocol around
the driver enable}, not the volts.

\question{Q3.2}{Your RS422 link is solid at 1\,m on the bench but throws
framing errors at 50\,m in the rig. Walk the diagnosis.}
Distance changes three things: signal integrity (reflections from missing/
incorrect termination, now significant at 50\,m), accumulated baud/clock skew
(less likely if it worked at all, but check both ends' crystals), and noise
pickup/ground shift (less for differential, but a floating or hugely shifted
common-mode can exceed the receiver's range). Method: scope the differential
pair \emph{at the receiver} --- ringing/overshoot says termination
(add/correct 120\,$\Omega$); a clean but late/early sample says baud; a
collapsed differential says a broken pair or common-mode out of range (check the
ground/shield and the receiver's common-mode limits). \trap the junior reflex is
``lower the baud'' --- which can mask a termination fault rather than fix it.
The senior method localises it on the scope at the receiver and fixes the
actual cause; it's the oscilloscope-based root-cause work on the resume.

\question{Q3.3}{Parity ``passed'' but the data is still wrong. How is that
possible, and what do real protocols use instead?}
Parity only detects an \emph{odd} number of bit flips in a single character ---
two flips (or any even number) pass silently, and it gives no protection across
a multi-byte message or against dropped/duplicated bytes. So a frame can pass
every byte's parity yet be corrupt at the message level. Real protocols add a
\textbf{checksum or CRC} over the whole frame (a CRC catches all single/double-
bit and most burst errors), plus length and sequence fields. \trap candidates
over-trust parity; the depth is knowing its exact guarantee (odd-error,
per-character) and that telemetry/CCDL/1553B-style links use a frame CRC ---
which is why the resume's packet parsing validates a frame checksum, not just
UART parity.

\question{Q3.4}{Two devices ``don't talk'' over RS232/RS422 even though both
are configured for 115200 8N1. Enumerate the likely physical-layer faults.}
For RS232: TX/RX not crossed (both expecting the other to be the DCE/DTE ---
need a null-modem swap), common ground missing, or inverted polarity. For
RS422/485: the differential pair A/B swapped (inverts the data $\rightarrow$
constant framing errors), TX pair of one wired to TX of the other instead of to
RX, missing termination, or (485) both ends stuck driving. Plus the shared
suspects: baud/format actually mismatched despite the label, flow-control lines
not satisfied (CTS held off so nothing sends), and one side's driver not
enabled. \trap the most common and most overlooked is \emph{A/B (or TX/RX)
swapped} --- it looks like total silence or all-framing-errors, and the
disciplined check is a scope/meter on the pair plus confirming the
crossover, before touching software.

\question{Q3.5}{Design an RS422 telemetry receiver that won't lose data and can
recover from a mid-frame glitch. What are the building blocks?}
Layered: (1) hardware --- correct termination and a crystal clock so framing is
reliable; (2) UART RX via DMA into a circular buffer (Chapter 4) so no byte is
lost even at high rate, with the IDLE-line interrupt marking frame boundaries;
(3) a framing protocol with a sync/start pattern, a length field, and a
\textbf{CRC} so a corrupted or mis-aligned frame is detected and discarded;
(4) a resynchronisation rule --- on CRC failure, scan forward for the next sync
pattern rather than trusting byte alignment. \trap the trap is assuming clean
byte framing implies clean \emph{message} framing; noise can insert/drop a byte
and shift every subsequent field. The robust receiver never trusts alignment ---
it re-finds the sync word and validates the CRC every frame, exactly the
discipline behind the resume's RS422/telemetry packet validation.

%=====================================================================
\section{Coding Gym}
%=====================================================================

\begin{partnote}
All ten problems are host-compilable and verified under
\code{gcc -Wall -Wextra -std=c11}. They are the arithmetic and bit-level
operations you actually perform when bringing up and validating a UART/RS422
link.
\end{partnote}

%---------------------------------------------------------------
\begin{gymproblem}{UART baud divider and error}

\textbf{Problem.} For a peripheral clock and target baud (16$\times$
oversampling), compute the \code{BRR} register, the actual baud, and the error
in ppm.

\textbf{Hints.}
\begin{itemize}
\item \code{BRR = round(f\_ck / baud)}; actual baud \code{= f\_ck / BRR}.
\item Error in ppm \code{= (actual - baud) * 1e6 / baud} (use 64-bit).
\end{itemize}

\attempt

\textbf{Solution.}
\begin{cgym}
#include <stdint.h>
#include <stdio.h>

static uint32_t baud_brr(uint32_t fck, uint32_t baud) {
    return (fck + baud / 2u) / baud;          /* rounded divisor */
}
static uint32_t baud_actual(uint32_t fck, uint32_t brr) {
    return fck / brr;
}
static int32_t baud_err_ppm(uint32_t actual, uint32_t baud) {
    return (int32_t)(((int64_t)actual - (int64_t)baud) * 1000000 / baud);
}

int main(void) {
    uint32_t brr = baud_brr(16000000u, 115200u);     /* 139 */
    uint32_t act = baud_actual(16000000u, brr);       /* 115107 */
    printf("brr=%u actual=%u err=%dppm\n", brr, act, baud_err_ppm(act, 115200u));
    /* brr=139 actual=115107 err=-807ppm */
    return 0;
}
\end{cgym}

$16{,}000{,}000/115200 = 138.9 \rightarrow 139$, giving $115{,}107$ baud, only
$-807$\,ppm ($-0.08\%$) off --- comfortably within tolerance. The error is worst
at high baud (small divisor), which is exactly where bring-up bugs appear, so
this calculation is the first thing you check when a fast link misbehaves.

\followups
\begin{enumerate}
\item \emph{``Why does 14.7456\,MHz give zero error?''} $\rightarrow$ It's an
  integer multiple of common bauds, so the divisor is exact.
\item \emph{``Add the OVER8 (8$\times$) mode.''} $\rightarrow$ The divider and
  the fractional field change; recompute with the 8$\times$ formula.
\item \emph{``At what baud does a 1\,MHz clock break 115200?''} $\rightarrow$
  Divisor 8.68 rounds to 9 $\rightarrow$ 111111 baud, $\sim$3.5\% error: too
  much.
\end{enumerate}
\end{gymproblem}

%---------------------------------------------------------------
\begin{gymproblem}{Frame (byte) transmission time}

\textbf{Problem.} Compute how long one character takes on the wire given baud
and format.

\textbf{Hints.}
\begin{itemize}
\item bits/char = 1 start + data + parity + stop.
\item time = bits / baud.
\end{itemize}

\attempt

\textbf{Solution.}
\begin{cgym}
#include <stdint.h>
#include <stdio.h>

static uint32_t frame_time_ns(uint32_t baud, int databits, int parity, int stopbits) {
    int bits = 1 + databits + (parity ? 1 : 0) + stopbits;
    return (uint32_t)((uint64_t)bits * 1000000000u / baud);
}

int main(void) {
    printf("frame=%u ns\n", frame_time_ns(115200u, 8, 0, 1));   /* 8N1 -> ~86805 ns */
    return 0;
}
\end{cgym}

8N1 is 10 bit-times, so at 115200 baud a byte takes $10/115200 \approx
86.8\,\mu$s, and the maximum throughput is $115200/10 = 11520$ bytes/s. This is
the number you use to size buffers and decide between interrupt and DMA: at
86.8\,$\mu$s/byte an ISR per byte is fine; at 1\,Mbaud (10\,$\mu$s/byte) it
isn't.

\followups
\begin{enumerate}
\item \emph{``Throughput in bytes/s for 8E2?''} $\rightarrow$ 11 bits/byte
  $\rightarrow$ baud/11.
\item \emph{``How long to send a 256-byte frame?''} $\rightarrow$
  256 $\times$ frame\_time; the basis of a transmit-timeout.
\item \emph{``Why is the IDLE detection time one frame?''} $\rightarrow$ A gap
  longer than a frame time signals end-of-message.
\end{enumerate}
\end{gymproblem}

%---------------------------------------------------------------
\begin{gymproblem}{Compute a parity bit}

\textbf{Problem.} Return the even-parity bit of a byte without a loop over all
8 bits being the only option --- use an XOR fold.

\textbf{Hints.}
\begin{itemize}
\item Fold the byte in halves with XOR: \code{b\^{}=b>>4; b\^{}=b>>2;
  b\^{}=b>>1}.
\item The low bit is 1 iff the popcount is odd $\rightarrow$ the even-parity
  bit.
\end{itemize}

\attempt

\textbf{Solution.}
\begin{cgym}
#include <stdint.h>
#include <stdio.h>

static int parity_even(uint8_t b) {
    b ^= b >> 4;
    b ^= b >> 2;
    b ^= b >> 1;
    return b & 1;          /* 1 if an odd number of bits were set */
}

int main(void) {
    printf("p(0x07)=%d p(0x03)=%d\n", parity_even(0x07), parity_even(0x03));
    /* 0x07 has 3 ones -> 1 ; 0x03 has 2 ones -> 0 */
    return 0;
}
\end{cgym}

The XOR fold computes parity in $O(\log n)$ and branch-free. The returned bit is
the \emph{even}-parity bit: appended, it makes the total number of 1s even.
Odd parity is its complement. Recompute on receive and compare to the received
parity bit; remember the guarantee is only ``odd number of flips, single
character.''

\followups
\begin{enumerate}
\item \emph{``Odd parity?''} $\rightarrow$ \code{return (b \& 1) \^{} 1;}.
\item \emph{``Parity of a 32-bit word?''} $\rightarrow$ Extend the fold:
  \code{x\^{}=x>>16}, then the same steps.
\item \emph{``Hardware does parity --- why compute it?''} $\rightarrow$ For a
  9-bit/manual scheme, a software protocol, or to verify the peripheral.
\end{enumerate}
\end{gymproblem}

%---------------------------------------------------------------
\begin{gymproblem}{Build a UART frame bit-sequence}

\textbf{Problem.} Produce the on-wire bit sequence (start, data LSB-first,
optional parity, stop) for a byte --- the basis of a software (bit-banged) UART
or a test vector generator.

\textbf{Hints.}
\begin{itemize}
\item Start = 0, data LSB first, parity if enabled, stop = 1.
\end{itemize}

\attempt

\textbf{Solution.}
\begin{cgym}
#include <stdint.h>
#include <stdio.h>

/* pmode: 0=none, 1=even, 2=odd. Returns number of bits written. */
static int build_frame(uint8_t data, int pmode, uint8_t *bits) {
    int n = 0, ones = 0;
    bits[n++] = 0;                                  /* start bit */
    for (int i = 0; i < 8; i++) {                   /* data, LSB first */
        int b = (data >> i) & 1;
        bits[n++] = (uint8_t)b;
        ones += b;
    }
    if (pmode) {
        int p = (pmode == 1) ? (ones & 1) : ((ones & 1) ^ 1);  /* even / odd */
        bits[n++] = (uint8_t)p;
    }
    bits[n++] = 1;                                  /* stop bit */
    return n;
}

int main(void) {
    uint8_t bits[12];
    int n = build_frame(0x41, 1, bits);             /* 'A', even parity */
    printf("nbits=%d start=%d stop=%d\n", n, bits[0], bits[n - 1]);  /* 11 0 1 */
    return 0;
}
\end{cgym}

This is literally what a bit-banged UART shifts out (one bit per timer tick at
the baud period) and what you feed a protocol-analyser model to validate a real
UART's output. Note \textbf{LSB-first} --- get the bit order backwards and every
byte is reversed, a classic ``the data is garbage but framing is fine'' bug.

\followups
\begin{enumerate}
\item \emph{``Bit-bang it: how often do you output each bit?''} $\rightarrow$
  Every baud period (\code{1/baud}) from a timer ISR.
\item \emph{``Add a second stop bit.''} $\rightarrow$ Append another \code{1};
  the receiver just sees a longer idle.
\item \emph{``9-bit data?''} $\rightarrow$ Loop to 9; the 9th bit often carries
  address/parity (next problem).
\end{enumerate}
\end{gymproblem}

%---------------------------------------------------------------
\begin{gymproblem}{Detect a framing error}

\textbf{Problem.} Given a sampled frame's bits, decide whether framing is valid
(start low, stop high).

\textbf{Hints.}
\begin{itemize}
\item Only the first and last sampled bits define framing validity here.
\end{itemize}

\attempt

\textbf{Solution.}
\begin{cgym}
#include <stdint.h>
#include <stdio.h>

static int framing_ok(const uint8_t *bits, int n) {
    return bits[0] == 0 && bits[n - 1] == 1;   /* start must be 0, stop must be 1 */
}

int main(void) {
    uint8_t good[] = {0, 1,0,0,0,0,0,1,0, 1};   /* start=0 ... stop=1 */
    uint8_t bad[]  = {0, 1,0,0,0,0,0,1,0, 0};   /* stop sampled as 0  */
    printf("good=%d bad=%d\n", framing_ok(good, 10), framing_ok(bad, 10));
    return 0;
}
\end{cgym}

A real UART raises the framing-error flag when it samples the stop position low.
The dominant cause is baud mismatch (the sample point has drifted onto a data
bit), so a flood of framing errors on bring-up points at the clock/baud config,
not the cable. Noise and break conditions are the other causes.

\followups
\begin{enumerate}
\item \emph{``How does the UART know where the stop bit is?''} $\rightarrow$ It
  counts bit-times from the start edge at the configured baud.
\item \emph{``Distinguish a break from a framing error.''} $\rightarrow$ A
  break holds the line low for a whole frame (all zeros incl.\ stop).
\item \emph{``Why recover by hunting the next idle?''} $\rightarrow$ After a
  framing error, byte alignment is suspect; resync on idle/start.
\end{enumerate}
\end{gymproblem}

%---------------------------------------------------------------
\begin{gymproblem}{RS485 9-bit address-mark frame}

\textbf{Problem.} Build a 9-bit frame that marks a byte as an address (9th bit
set) and a test that recognises address bytes.

\textbf{Hints.}
\begin{itemize}
\item The 9th bit (bit 8) distinguishes address from data.
\end{itemize}

\attempt

\textbf{Solution.}
\begin{cgym}
#include <stdint.h>
#include <stdio.h>

static uint16_t make_addr_frame(uint8_t a) { return (uint16_t)(a | 0x100u); }  /* set bit 8 */
static int      is_address(uint16_t f)      { return (f >> 8) & 1; }

int main(void) {
    uint16_t af = make_addr_frame(0x12);
    printf("frame=0x%03X is_addr=%d data_is_addr=%d\n",
           af, is_address(af), is_address(0x0041));   /* 0x112 1 0 */
    return 0;
}
\end{cgym}

In multidrop RS485 the master sends the slave's address with bit 8 set; slaves
in address-mute mode wake only on an address byte, compare it, and the matching
slave then accepts the following data bytes (bit 8 clear). The hardware
``mute until addressed'' feature means non-addressed nodes aren't interrupted on
every byte --- essential on a busy shared bus.

\followups
\begin{enumerate}
\item \emph{``How does a slave stop listening after its message?''}
  $\rightarrow$ Re-enter mute on the next address byte or a frame terminator.
\item \emph{``Collisions on the bus?''} $\rightarrow$ Only the master initiates;
  slaves answer only when polled --- the protocol prevents contention.
\item \emph{``9-bit vs parity reuse.''} $\rightarrow$ Some UARTs repurpose the
  parity-bit slot as the 9th data/address bit.
\end{enumerate}
\end{gymproblem}

%---------------------------------------------------------------
\begin{gymproblem}{RS485 driver-enable guard time}

\textbf{Problem.} Compute how long (ns) to hold the DE line after the last bit
before releasing the bus, expressed as a number of bit-periods.

\textbf{Hints.}
\begin{itemize}
\item Guard time = guard\_bits / baud.
\end{itemize}

\attempt

\textbf{Solution.}
\begin{cgym}
#include <stdint.h>
#include <stdio.h>

static uint32_t de_guard_ns(uint32_t baud, int guard_bits) {
    return (uint32_t)((uint64_t)guard_bits * 1000000000u / baud);
}

int main(void) {
    printf("guard=%u ns\n", de_guard_ns(115200u, 2));   /* ~17361 ns for 2 bits */
    return 0;
}
\end{cgym}

You assert DE before the first start bit and hold it for the whole frame plus a
guard of a bit or two after the final stop bit, so the last byte fully shifts
out of the transceiver before you tri-state --- release too early and the last
character is clipped. At 115200 baud a 2-bit guard is $\sim$17.4\,$\mu$s. Many
UARTs automate this with a hardware DE signal and a programmable
assertion/de-assertion time, which removes the software-timing race entirely.

\followups
\begin{enumerate}
\item \emph{``Why not release DE immediately on TX-complete?''} $\rightarrow$
  TXE (register empty) $\neq$ TC (shift register empty); wait for TC plus guard.
\item \emph{``Turnaround budget on a polled bus?''} $\rightarrow$ Master must
  wait DE-release + cable delay before expecting the slave's reply.
\item \emph{``Hardware DE control?''} $\rightarrow$ Configure the UART's
  driver-enable with assertion/deassertion times; no ISR needed.
\end{enumerate}
\end{gymproblem}

%---------------------------------------------------------------
\begin{gymproblem}{Majority-vote bit sampling}

\textbf{Problem.} A receiver oversamples each bit and takes three samples near
the centre; return the de-noised bit by majority vote.

\textbf{Hints.}
\begin{itemize}
\item Majority of three = at least two agree.
\end{itemize}

\attempt

\textbf{Solution.}
\begin{cgym}
#include <stdio.h>

static int majority3(int a, int b, int c) {
    return (a + b + c) >= 2;     /* at least two of three are 1 */
}

int main(void) {
    printf("maj(1,0,1)=%d maj(0,0,1)=%d\n", majority3(1, 0, 1), majority3(0, 0, 1));
    return 0;
}
\end{cgym}

Real UART receivers oversample 16$\times$ and sample bits 8, 9, 10 (around the
centre), taking the majority to reject a single noise glitch on the line ---
this is why an oversampling UART is more robust than a naive ``sample once at
the middle'' bit-bang. The same idea hardens a software UART or a noisy GPIO
input.

\followups
\begin{enumerate}
\item \emph{``Why sample at the centre, not the edge?''} $\rightarrow$ The
  signal is most settled mid-bit; edges are where transitions/ringing live.
\item \emph{``5-sample majority?''} $\rightarrow$ \code{sum >= 3}; more noise
  immunity, more compute.
\item \emph{``How does oversampling set baud tolerance?''} $\rightarrow$ More
  samples per bit $\rightarrow$ finer sample-point placement $\rightarrow$
  better drift tolerance.
\end{enumerate}
\end{gymproblem}

%---------------------------------------------------------------
\begin{gymproblem}{Frame checksum (LIN-style, with carry)}

\textbf{Problem.} Compute a single-byte checksum over a frame using a
sum-with-carry then one's-complement (the LIN ``classic'' checksum).

\textbf{Hints.}
\begin{itemize}
\item Add each byte; when the sum exceeds 8 bits, fold the carry back in.
\item Complement the final sum.
\end{itemize}

\attempt

\textbf{Solution.}
\begin{cgym}
#include <stdint.h>
#include <stdio.h>

static uint8_t lin_checksum(const uint8_t *d, int n) {
    uint16_t s = 0;
    for (int i = 0; i < n; i++) {
        s = (uint16_t)(s + d[i]);
        if (s > 0xFFu) s = (uint16_t)((s & 0xFFu) + 1u);   /* fold the carry */
    }
    return (uint8_t)(~s);                                  /* one's complement */
}

int main(void) {
    uint8_t d[] = {0x55, 0x93, 0xE5};
    printf("cksum=0x%02X\n", lin_checksum(d, 3));   /* 0x31 */
    return 0;
}
\end{cgym}

A frame checksum catches errors a per-character parity bit misses --- dropped
bytes, two-bit flips, mis-ordering. This carry-folded sum is cheap and used in
LIN and many simple telemetry framings; for stronger protection (burst errors)
you step up to a CRC (Chapter 13). The receiver recomputes over the received
bytes and compares; a mismatch discards the frame.

\followups
\begin{enumerate}
\item \emph{``Why fold the carry instead of dropping it?''} $\rightarrow$ It
  makes the checksum sensitive to all byte values (a one's-complement sum),
  catching more errors.
\item \emph{``Checksum vs CRC --- when each?''} $\rightarrow$ Checksum for cheap
  low-rate links; CRC for noisy/critical links (catches burst errors).
\item \emph{``Include the header/length in the sum?''} $\rightarrow$ Yes ---
  cover every field you want protected, consistently on both ends.
\end{enumerate}
\end{gymproblem}

%---------------------------------------------------------------
\begin{gymproblem}{Baud-mismatch tolerance check}

\textbf{Problem.} Given a baud error in ppm and the frame length in bits,
decide whether the cumulative sampling error stays within half a bit.

\textbf{Hints.}
\begin{itemize}
\item Cumulative error $\approx$ err\_ppm $\times$ bits; half a bit is
  500000\,ppm.
\end{itemize}

\attempt

\textbf{Solution.}
\begin{cgym}
#include <stdint.h>
#include <stdio.h>

static int baud_match_ok(int32_t err_ppm, int bits_per_frame) {
    long cum = (long)err_ppm * bits_per_frame;
    if (cum < 0) cum = -cum;
    return cum < 500000;          /* < half a bit (500000 ppm) */
}

int main(void) {
    printf("ok(-807,10)=%d ok(60000,10)=%d\n",
           baud_match_ok(-807, 10), baud_match_ok(60000, 10));
    /* -807ppm*10 = 8070 -> ok ; 60000ppm*10 = 600000 -> fail */
    return 0;
}
\end{cgym}

This formalises ``the sample point must not drift past the bit edge.'' A
$-807$\,ppm error over a 10-bit frame is 8070\,ppm of cumulative drift ---
nowhere near the half-bit (500000\,ppm) limit --- so it works. A 6\% (60000\,ppm)
error fails: by the stop bit you're sampling the wrong bit. Both ends' errors
add, so budget half the limit to each side.

\followups
\begin{enumerate}
\item \emph{``Two stop bits --- does tolerance change?''} $\rightarrow$ Slightly
  more bits to the last sample, marginally tighter; the start-bit resync is
  what bounds it per character.
\item \emph{``Why does each character resync help?''} $\rightarrow$ The start
  edge re-zeroes the timing every byte, so error doesn't accumulate across
  bytes.
\item \emph{``Budget split between two crystals.''} $\rightarrow$ Give each end
  $<$1\% so the sum stays safe.
\end{enumerate}
\end{gymproblem}

%=====================================================================
\section{Link to Her Resume}
%=====================================================================

\begin{resumelink}
\textbf{RS422 daily work:} ``RS422 is my daily interface on the DFCC --- I
develop and validate the Embedded C routines for card-to-card RS422 links.
I can explain why it's RS422 and not RS232: differential pairs reject the
common-mode noise of an EMI-heavy aircraft (MIL-STD-461), run long and fast,
and separate command and telemetry full-duplex.''

\medskip
\textbf{Bring-up diagnosis:} ``When a differential link fails at distance I
scope the pair \emph{at the receiver} --- ringing means termination, a drifted
sample point means baud/clock, a swapped A/B means inverted data. I fix the
actual cause rather than just dropping the baud.''

\medskip
\textbf{Frame integrity:} ``Parity only catches odd single-character errors, so
the telemetry framing I validate uses a frame checksum/CRC plus a sync word and
length, and re-finds the sync on any failure rather than trusting byte
alignment.''

\medskip
\small Maps to concrete resume lines: RS422 Embedded C interface routines on the
DFCC at TASL, oscilloscope-based root-cause analysis, and RS422 telemetry packet
validation. Keep claims to what the resume states; the deep project scripting is
Chapter 15, from \code{inputs/INTAKE.md} only.
\end{resumelink}

%=====================================================================
\section{Self-Test Scorecard}
%=====================================================================

\begin{scorecard}
\begin{enumerate}
\item UART vs RS232/422/485 --- which is the framing and which is the
  electrical standard?
\item Describe an 8N1 frame and state how many bit-times it is.
\item 16\,MHz clock, 115200 baud, 16$\times$ oversample --- give \code{BRR} and
  the approximate error.
\item Single-ended vs differential --- which rejects common-mode noise, and
  why?
\item RS422 vs RS485 --- the real distinguishing property (not the volts).
\item What exactly does a parity bit guarantee, and what does it miss?
\item Name the three classic UART receive errors.
\item Why must RS485 manage the driver-enable line, and what's the timing risk?
\item Where and why do you terminate a differential pair?
\item Why do flight computers use RS422 for inter-card links?
\end{enumerate}
\end{scorecard}

\begin{answerkey}
\begin{enumerate}
\item UART = the asynchronous framing (peripheral); RS232/422/485 = electrical
  standards (voltages, single-ended vs differential).
\item Start (0), 8 data LSB-first, no parity, 1 stop (1) --- 10 bit-times.
\item \code{BRR}=139, actual 115107 baud, $\approx-0.08\%$ ($-807$\,ppm).
\item Differential --- noise hits both wires equally (common-mode) and cancels
  at the difference receiver.
\item Topology/duplex: RS422 is one driver, usually full-duplex point-to-point;
  RS485 is multi-driver half-duplex needing DE arbitration.
\item It detects an odd number of bit flips in a single character; it misses
  even-count errors and any multi-byte/dropped-byte corruption.
\item Framing error (stop sampled low), parity error, overrun (unread byte
  overwritten).
\item It's half-duplex shared-bus, so each driver must tri-state except when
  transmitting; the risk is turnaround --- releasing DE too early clips the last
  byte, too late collides.
\item At each end of the pair, with the cable's characteristic impedance
  ($\sim$120\,$\Omega$), to absorb reflections that would otherwise corrupt fast
  edges.
\item Differential noise immunity (EMI/MIL-STD-461), long reliable range,
  higher baud, and full-duplex separation of command/telemetry --- vs RS232's
  short single-ended reach.
\end{enumerate}
\end{answerkey}

\begin{partnote}
\emph{End of Chapter 5.} Next: \textbf{SPI \& I2C} --- the two short-range
synchronous buses: clock polarity/phase modes, chip-select and addressing,
clock stretching and arbitration, and when to choose each.
\end{partnote}
